\documentclass[12pt, titlepage]{article}
\usepackage{hyperref}

\title{Feedback Controls Lab Weekly Report \\ \normalsize{Week 2}}
\author{Aydan Vissing, Brady Schick, Gabriel Southwell, Simon Ross}
\date{\today}

\begin{document}
\maketitle

\section*{Aydan Vissing (6 hrs spent)}

I researched more on PID controls and what we would need for the drone project. We then had a meeting with Dr. Fredette to discuss what the proposal would look like. I then participated in the Circuit board lab, where we learned about what we will need if we decide to make our own PCB for the drone kit. I then did some more research into what we would need for the proposal and what we could do to stay under budget since it dropped a ton. I was able to determine that as long as we only buy one or two premade test kits or drones, we would have enough to make several prototypes for ourselves. This could be revised if we happen to find some cheaper options, but the average cost for what seems to be a reliable kit is around \$200.

\section*{Brady Schick (7.5 hrs spent)}

On Tuesday this week, our group spent 2.5 hours going through the old controls labs in order to get a better understanding of what we need to accomplish with our project. On Wednesday, I also spent 30 minutes listing some marketing requirements that will need to be met. On Thursday, we spent two hours on the PCB lab since that is something that will be useful if we decide we need to custom-design our drones for the lab. On Friday, we met for an hour to develop a more official list of marketing requirements. On Saturday, we completed the old controls labs from Tuesday.

\section*{Gabriel Southwell (8 hrs spent)}

For this week, I did a few things to broaden my understanding of control systems and gain a better understanding of what we are trying to teach in our lab. On Tuesday, some of us started to work on the old labs to see what they were trying to teach. We finished the labs on Saturday.

We all also went to Dr. Kohl’s PCB design lab. 

On Friday we met with Dr. Fredette and talked about our project’s marketing requirements. This also gave us a chance to talk more about the goals for the drones and clear some of my confusion up. 

I researched a bunch about drone firmware. Right now I see two good options:
\begin{itemize}
	\item Start with Betaflight and strip it down to meet our needs 
	\item Start from the bottom and copy betaflight’s homework
    \begin{itemize}
		\item It would be based on freeRTOS since that already supports basically every microcontroller
	\end{itemize}
\end{itemize}

\section*{Simon Ross (6.5 hrs spent)}

This week, our team worked on thinking about our proposal, developing the skills we’ll need for the rest of the project, and researching viable options going forward. Ultimately, I think the research completed by Gabe will be particularly valuable as we seem to be headed toward creating our own drone--I personally need to buckle down on the research side and see what else we need to know.

The main things I accomplished this week were researching PID controllers, completing some of the old controls labs, and attending the PCB workshop with Dr. Kohl. Overall, I have greatly increased my understanding of basic controllers and will be able to think about the requirements for our project more clearly. Working through the old labs has also helped with this, as it allowed me to practically see analog controllers working via physical feedback loops. Hopefully our team will be able to create a similar “mind blown” experience in our project to the experience I had over the labs. Finally, the PCB workshop seems like invaluable experience if we need to create our own PCB’s for the project (which is looking more likely). 

The project is in a good place–-we discussed marketing requirements on Fredette Friday, which will allow us to complete our proposal in the upcoming days. Aydan’s fears appear to be founded, as we haven’t really discovered a long-standing kit that supplies all of the components we need. I’m a proponent of finding as many open-source components (software and hardware) as possible so that we don’t need to do everything from scratch. We will continue to discuss things as a team. For now, to sum things up, this next week will be wrapping up the old labs, working on the project proposal, and researching what all we will need to create a functional drone that can accomplish our goals. 


\end{document}